\documentclass[10pt,letterpaper]{article}
\usepackage[margin=0.65in,headheight=14pt]{geometry}
\usepackage[T1]{fontenc}
\usepackage{lmodern,amsmath,amssymb,booktabs,graphicx,xcolor,microtype,fancyhdr}
\usepackage[colorlinks=true,urlcolor=teal]{hyperref}
\definecolor{navy}{HTML}{192B4D}
\definecolor{teal}{HTML}{007D80}
\definecolor{pale}{HTML}{F3F6F8}
\pagestyle{fancy}\fancyhf{}
\fancyhead[L]{\small MF 731 | Homework 3}
\fancyhead[R]{\small VaR Backtesting}
\fancyfoot[C]{\small\thepage}
\setlength{\parindent}{0pt}
\setlength{\parskip}{4pt}
\hypersetup{pdftitle={MF 731 Homework 3 - VaR Backtesting}}
\input{var_backtest_numbers.tex}
\begin{document}
{\LARGE\bfseries\textcolor{navy}{One-Day VaR Backtesting}}

\textbf{1. Data.} Use the S\&P 500 daily closing index from Yahoo Finance
(\texttt{\^{}GSPC}), from 2016-05-23 through 2026-05-22, inclusive.
There are $N=\PriceCount$ prices and \ReturnCount{} log returns.
The first and last five observations are:
\begin{center}
\begin{tabular}{lrrr}\toprule
Date & Close & Date & Close\\\midrule
\DataRowA
\DataRowB
\DataRowC
\DataRowD
\DataRowE
\bottomrule\end{tabular}
\end{center}

\textbf{2. Forecasts.} For $x_t=\log(P_t/P_{t-1})$, the loss on the
constant \$1 position is $L_t=1-e^{x_t}$. Each forecast uses only returns
known before day $t$. Historical VaR is the 95th percentile of the full
losses from the previous $M=1010$ returns, using the empirical inverse
CDF (the 960th smallest loss).

Initialize the normal model with the first 1010 returns:
$\widehat\mu=\InitialMean$, $\widehat\sigma=\InitialSD$
(sample standard deviation). After each observed return, update
\[
\widehat\mu_{t+1}=0.97\widehat\mu_t+0.03x_t,\qquad
\widehat\sigma_{t+1}^2=0.97\widehat\sigma_t^2
 +0.03(x_t-\widehat\mu_t)^2.
\]
The day-$t$ normal VaR is
$1-\exp\!\big(\widehat\mu_t+\widehat\sigma_t\Phi^{-1}(0.05)\big)$.
Both methods evaluate the same $J=N-1-M=\TestCount$ returns,
from \TestStart{} to 2026-05-22. Since $N$ counts prices, differencing
loses one observation; this gives $J=N-1-M$ rather than the handout's $N-M$.

\begin{center}
\includegraphics[width=\linewidth]{var_backtest_plot.png}
\end{center}

\textbf{3. Exceedance counts.} Count an exceedance when $L_t>\mathrm{VaR}_t$.
Under correct coverage with independent exceedance indicators,
$K\sim\mathrm{Binomial}(\TestCount,0.05)$, so $E[K]=\ExpectedCount$.
The central 95\% binomial interval is
$[F_K^{-1}(0.025),F_K^{-1}(0.975)]=[\CountLower,\CountUpper]$.
\begin{center}
\begin{tabular}{lrrc}\toprule
Method & Exceedances & Rate & 95\% count interval\\\midrule
Historical & \HistCount & \HistRate & \HistDecision\\
Normal EWMA & \NormalCount & \NormalRate & \NormalDecision\\
\bottomrule\end{tabular}
\end{center}
\CountConclusion

\begin{center}
\colorbox{pale}{\parbox{0.94\linewidth}{\centering
\textbf{Python code and running instructions}\\[3pt]
\url{https://github.com/jesuswalking/MF-731-HW1/tree/main/HW3}}}
\end{center}
\end{document}
